\documentclass[10pt]{article}
\usepackage[spanish]{babel}
\usepackage{cv}
\renewcommand{\techlabel}{Tecnologías}
\hypersetup{pdftitle={CV - Pedro Gonzalez Nuñez}, pdfauthor={Pedro Gonzalez Nuñez}}

\begin{document}

\cvheader{photo.jpg}{Pedro Gonzalez Nuñez}
  {7 de septiembre de 2002}{pedrogonzalezhudson@gmail.com}
  {+54 9 11 5029 7907}{Buenos Aires, Argentina}
  {www.linkedin.com/in/pedrogonzaleznunez/}{github.com/pedrogonzaleznunez}

\vspace{4pt}
\cvsection{Sobre mí}{%
  Estudiante de cuarto año de Ingeniería Informática con una sólida base en desarrollo de software,
  automatización y sistemas de inteligencia artificial. Tengo experiencia tanto en proyectos de
  investigación académica como en entornos de desarrollo profesional, combinando programación,
  análisis matemático y trabajo en equipo. Me interesa especialmente aplicar la tecnología para
  resolver problemas complejos y optimizar procesos en contextos reales y multidisciplinarios.%
}

\sectionrule
% ---------- EXPERIENCE: mismo orden que LinkedIn; cada \jobrow = fila de dos trabajos alineados ----------
\explabel{Experiencia Laboral}
\jobrow{%
  \job{Software Engineer}{Galo AI | Medio tiempo | Agosto 2026 -- Septiembre 2026}
      {Prometheus, Loki, Grafana, Metabase, Cloudflare, VPS, DNS, Cloud Storage.}{
      \item Migré la infraestructura de un único servidor a una arquitectura distribuida de tres servidores, aislando el monitoreo (Prometheus, Loki, Metabase) para eliminar puntos únicos de falla.
      \item Armé un pipeline de despliegue con staging previo y dashboards de Grafana para observabilidad en tiempo real.
      \item Aseguré el almacenamiento cloud con signed URLs e interconecté servicios entre VPS con Cloudflare Tunnels (Zero Trust).
    }%
}{%
  \job{Forward Deployed Engineer}{Galo AI | Tiempo completo | Abril 2026 -- Julio 2026}
      {Python, SQL, OpenAI/Anthropic APIs, RAG, Prompt Engineering, LLM Ops.}{
      \item Desarrollé y desplegué workflows de agentes de IA en producción con validaciones, reintentos y guardrails.
      \item Implementé sistemas RAG y structured outputs con tool calling y JSON schemas dinámicos.
      \item Armé pipelines de evaluación y regresión de LLMs para reducir alucinaciones.
    }%
}
\jobrow{%
  \job{Software Engineer}{Codify Company | Medio tiempo | Septiembre 2025 -- Febrero 2026}
      {React, Next.js, NestJS, Prisma, PostgreSQL, Docker, Coolify, AWS.}{
      \item Desarrollé features full stack para una plataforma de alta concurrencia con muchos usuarios activos.
      \item Resolví incidentes en producción y optimicé endpoints críticos, mejorando los tiempos de respuesta.
    }%
}{%
  \job{Full-stack Developer}{ITBA Rocketry Team | Contrato | Agosto 2025 -- Septiembre 2025}
      {Next.js, React.}{
      \item Desarrollé herramientas internas y aplicaciones web para el equipo, con foco en frontend.
      \item Desarrollé la landing page del equipo.
    }%
}
\jobrow{%
  \job{Research Assistant}{ITBA | Medio tiempo | Octubre 2024 -- Septiembre 2025}
      {Python, Signal Processing.}{
      \item Desarrollé un robot integrando electrónica y sistemas de control.
      \item Implementé su parte electrónica: batería, motores y sistema de control.
    }%
}{%
  \job{AI Software Engineer}{Push\&Pop Growth | Freelance | Marzo 2025 -- Agosto 2025}
      {Python, FastAPI, Google ADK, PostgreSQL, WhatsApp API, Telegram API, Azure.}{
      \item Desarrollé agentes de IA con Google ADK, integrados con WhatsApp y Telegram.
      \item Desarrollé servicios REST con FastAPI y PostgreSQL, desplegados en Azure.
    }%
}

\sectionrule
\cvsection{Formación\\Académica}{%
  \edu{Instituto Tecnológico de Buenos Aires (ITBA)}{Ingeniería Informática | 2022 -- Actualidad}
  \edu{Cursos de Desarrollo de Software (Udemy)}{NestJS, Next.js, React, Docker | 2025 -- Actualidad}
  \edu{E.E.S.T. N.º~7 ``T.R.Q.'' (IMPA)}{Técnico en Aviónica | 2015 -- 2021}
  \edu{Curso de Negocios Digitales (Udemy)}{Ventas y manejo de clientes | Oct 2023 -- Dic 2023}%
}

\sectionrule
\cvsection{Intereses \&\\Hobbies}{%
  \inlinelist{Fútbol}{Trenes en miniatura}{Física aplicada (Relatividad y Electromagnetismo)}%
}

\sectionrule
\cvsection{Habilidades\\Técnicas}{%
  \textbf{Lenguajes:} TypeScript, Python, JavaScript, Java, C, C\#, HTML, CSS\\
  \textbf{Frontend \& Mobile:} React, Next.js, Desarrollo iOS, React Native\\
  \textbf{Backend:} Node.js (Express, NestJS), FastAPI, Django, Java Spring Boot\\
  \textbf{Bases de datos:} PostgreSQL, MySQL, MongoDB, SQLite3, Prisma (ORM)\\
  \textbf{Cloud, DevOps \& Herramientas:} AWS, Azure, Docker, Git, GitHub, Maven\\
  \textbf{Testing \& QA:} Testing automatizado (Unit, Integration, E2E), Jest, Cypress\\
  \textbf{Hardware:} Arduino, ESP32, Raspberry Pi%
}

\sectionrule
\cvsection{Idiomas}{%
  \lang{Español}{C2 (Nativo)}\lang{Inglés}{C1 (Avanzado)}\lang{Alemán}{B1 (Intermedio)}\lang{Francés}{A1 (Básico)}%
}

\end{document}
